% Preámbulo de la memoria.

%%%%%%%%%%%%%%%%%%%%%%%%%%%% Página, letra e idioma %%%%%%%%%%%%%%%%%%%%%%%%%%%%

\usepackage[a4paper,left=3cm,right=2.5cm,top=2.7cm,bottom=2.7cm,headheight=15pt]{geometry}
\usepackage[T1]{fontenc}
\usepackage{lmodern}
\usepackage[spanish,es-tabla,es-nosectiondot,es-noshorthands]{babel}
\usepackage{microtype}
\usepackage{csquotes}

% Interlineado de 1,2 sobre la letra de 11 pt (unos 16,3 pt entre líneas); párrafos sin sangría, separados por un espacio
\usepackage{setspace}
\setstretch{1.2}
\usepackage[skip=0.6\baselineskip plus 2pt,indent=0pt]{parskip}
\raggedbottom
\clubpenalty=10000
\widowpenalty=10000

% Capítulo: número grande en gris, título debajo alineado a la derecha y línea fina
\usepackage{xcolor}
\definecolor{numcap}{HTML}{B0B0B0}
\usepackage{titlesec}
\titleformat{\chapter}[display]
  {\normalfont\bfseries\raggedleft\hyphenpenalty=10000\exhyphenpenalty=10000}
  {\fontsize{90}{90}\selectfont\color{numcap}\thechapter}
  {-6pt}
  {\fontsize{32}{38}\selectfont}
  [{\vspace{4pt}\titlerule[0.8pt]}]
\titleformat{name=\chapter,numberless}[display]
  {\normalfont\bfseries\raggedleft\hyphenpenalty=10000\exhyphenpenalty=10000}
  {}
  {0pt}
  {\fontsize{32}{38}\selectfont}
  [{\vspace{4pt}\titlerule[0.8pt]}]
\titlespacing*{\chapter}{0pt}{-10pt}{36pt}
\titlespacing*{name=\chapter,numberless}{0pt}{40pt}{36pt}
\definecolor{azultitulo}{HTML}{1F2A44}
\definecolor{numcap}{HTML}{B8BCC6}
\titleformat{\chapter}[display]
  {\normalfont\bfseries\raggedleft\color{azultitulo}\hyphenpenalty=10000\exhyphenpenalty=10000}
  {\fontsize{90}{90}\selectfont\color{numcap}\thechapter}
  {-6pt}
  {\fontsize{32}{38}\selectfont}
  [{\vspace{4pt}\titlerule[0.8pt]}]
\titleformat{name=\chapter,numberless}[display]
  {\normalfont\bfseries\raggedleft\color{azultitulo}\hyphenpenalty=10000\exhyphenpenalty=10000}
  {}
  {0pt}
  {\fontsize{32}{38}\selectfont}
  [{\vspace{4pt}\titlerule[0.8pt]}]
\titleformat{\section}{\normalfont\Large\bfseries\color{azultitulo}}{\thesection}{0.8em}{}
\titleformat{\subsection}{\normalfont\large\bfseries\color{azultitulo}}{\thesubsection}{0.8em}{}

% Las listas de los índices, a interlineado sencillo (el título no cambia)
\makeatletter
\let\starttocoriginal\@starttoc
\renewcommand{\@starttoc}[1]{\begingroup\setstretch{1}\starttocoriginal{#1}\endgroup}
\makeatother

\addto\captionsspanish{%
  \renewcommand{\contentsname}{Índice}%
  \renewcommand{\appendixname}{Anexo}}

%%%%%%%%%%%%%%%%%%%%%%%%%%%% Figuras, tablas y números %%%%%%%%%%%%%%%%%%%%%%%%%%

\usepackage{graphicx}
\usepackage{xcolor}
\usepackage{amsmath,amssymb}
\usepackage{booktabs}
\usepackage{array}
\usepackage{multirow}
\usepackage{enumitem}
\usepackage{caption}
\captionsetup{font=small,labelfont=bf,skip=8pt}
\captionsetup[table]{position=top}
\usepackage{subcaption}
% [H] coloca una figura exactamente donde está en el texto, aunque deje hueco
% al final de la página: para las que se describen justo antes de aparecer
\usepackage{float}
% Ninguna figura cruza a la sección siguiente
\usepackage[section]{placeins}

% Cada figura y cada tabla dejan un pequeño margen propio por encima (14 pt), para
% que no queden pegadas al título o al texto que las precede, y su pie va junto a
% ellas: el espacio entre párrafos no se aplica dentro de figuras ni tablas.
% Figuras y tablas llevan el mismo margen: son el mismo tipo de elemento.
% El margen se añade al colocar la figura en medio del texto (en su sitio, o fija
% con [H]), no en el punto del código donde se escribe: si la figura se va a otra
% página, no deja un hueco en el texto. El \penalty0 deja cortar la página justo
% antes: sin él, una figura fija que no cabe obliga a partir por la mitad el
% párrafo que la precede.
% Las figuras pueden ocupar hasta el 90 % de lo alto de una página con texto
% debajo (por defecto, el 70 %): así dos figuras seguidas comparten página con el
% texto que las sigue, en lugar de quedar solas en una página de figuras
\renewcommand{\topfraction}{0.9}
\renewcommand{\textfraction}{0.1}
\renewcommand{\floatpagefraction}{0.8}
\setlength{\intextsep}{6pt plus 2pt minus 2pt}
\makeatletter
\patchcmd{\@addtocurcol}
  {\vskip\intextsep\box\@currbox}{\vskip\intextsep\vskip14pt\box\@currbox}
  {}{\PackageWarning{tfg}{No se pudo añadir el margen de las figuras en el texto}}
\patchcmd{\float@endH}
  {\@endfloatbox\vskip\intextsep}{\@endfloatbox\vspace*{14pt}\vskip\intextsep}
  {}{\PackageWarning{tfg}{No se pudo añadir el margen de las figuras fijas}}
\makeatother
\AtBeginEnvironment{figure}{\setlength{\parskip}{0pt}\par\penalty0}
\AtBeginEnvironment{table}{\setlength{\parskip}{0pt}\par\penalty0}

% Números según la RAE: coma decimal, espacio antes de las unidades y del %,
% y separación de miles solo a partir de cinco cifras (2377, pero 10 000)
\usepackage{siunitx}
\sisetup{
  output-decimal-marker = {,},
  group-separator = {\,},
  group-minimum-digits = 5,
}

% Diagramas (TikZ): cajas neutras y paneles de color suave por etapa. Cada
% diagrama, en figuras/
\usepackage{adjustbox} % reduce un diagrama que se pase del ancho del texto
\usepackage{tikz}
\usetikzlibrary{positioning,fit,arrows.meta,calc,backgrounds}
\definecolor{tinta}{HTML}{0b0b0b}
\definecolor{tinta2}{HTML}{52514e}
\definecolor{tenue}{HTML}{898781}
\definecolor{borde}{HTML}{b5b4aa}
\definecolor{fondo}{HTML}{f4f4f1}
\definecolor{panelaudio}{HTML}{eef3f9}
\definecolor{panelimagen}{HTML}{eef6f1}
\definecolor{panelgolpe}{HTML}{fbf4ea}
% Letra de las cajas: 9 pt fijos, el tamaño con el que se dimensionaron las cajas
% (\footnotesize crece con la letra del documento y las palabras largas se salen)
\newcommand{\letradiagrama}{\fontsize{9}{11}\selectfont}
\tikzset{
  caja/.style={draw=borde, fill=white, rounded corners=2pt, align=center,
    font=\letradiagrama, inner sep=3.5pt, minimum height=9mm, text width=21mm, text=tinta},
  entrada/.style={caja, fill=fondo},
  linea/.style={draw=tinta2, line width=0.55pt},
  flecha/.style={linea, -{Stealth[length=1.8mm]}},
  dato/.style={font=\scriptsize\itshape, text=tinta2, inner sep=1pt},
  panel/.style={rounded corners=6pt, inner sep=7pt},
  titulopanel/.style={font=\scriptsize, text=tinta2, anchor=south west, inner sep=2pt},
  seccion/.style={font=\tiny, text=tenue, anchor=north east, inner sep=1.5pt},
}
% Segunda línea, más pequeña y gris, dentro de una caja
\newcommand{\subcaja}[1]{\\[-1pt]{\scriptsize\color{tenue}#1}}
% Iconos de los diagramas: el jugador, un pictograma de remate como los de la
% figura de la ventana (0,66 de alto, pies en el origen), y la pelota de pádel
% (radio 0,085, centrada). El color se pasa al usarlos: {jugador=tinta}.
\definecolor{pelota}{HTML}{e2e23c}
\tikzset{
  pics/jugador/.style={code={
    \fill[#1] (0.015,0.44) circle (0.052);
    \draw[#1, line width=1.6pt, line cap=round] (0,0.35) -- (0,0.19);
    \draw[#1, line width=1.1pt, line cap=round, line join=round]
      (0,0.19) -- (-0.075,0.0) (0,0.19) -- (0.075,0.0)
      (0,0.33) -- (0.065,0.43) -- (0.085,0.52)
      (0,0.33) -- (-0.07,0.39) -- (-0.085,0.47);
    \draw[#1, line width=0.8pt, line cap=round] (0.085,0.52) -- (0.095,0.56);
    \fill[#1] (0.105,0.615) ellipse [x radius=0.032, y radius=0.045];
  }},
  pics/pelota/.style={code={
    \fill[#1, draw=tinta, line width=0.45pt] (0,0) circle (0.085);
    \draw[white, line width=0.55pt]
      (-0.07,0.035) .. controls (-0.02,-0.025) and (0.02,-0.025) .. (0.07,0.035);
  }},
}

%%%%%%%%%%%%%%%%%%%%%%%%%%%% Órdenes de terminal %%%%%%%%%%%%%%%%%%%%%%%%%%%%%%%%%

% Bloques con las órdenes que se escriben en la terminal: letra monoespaciada
% sobre el fondo gris de los diagramas, el «$» del indicador en gris y los
% comentarios en gris y cursiva. Uso: \begin{lstlisting}[style=terminal] ...
% \end{lstlisting} (el editor lo reconoce como código; con \begin{terminal} sale
% igual en el PDF, pero el editor colorea los «$» como si abrieran fórmulas)
\usepackage{listings}
\lstdefinestyle{terminal}{
  language=bash,
  basicstyle=\ttfamily\small\color{tinta},
  keywordstyle=\color{tinta},
  commentstyle=\color{tenue}\itshape,
  backgroundcolor=\color{fondo},
  frame=single, rulecolor=\color{borde}, framerule=0.4pt, framesep=5pt,
  xleftmargin=6pt, xrightmargin=6pt,
  aboveskip=10pt, belowskip=8pt,
  columns=fullflexible, keepspaces=true, showstringspaces=false,
  breaklines=true,
  literate={\$}{{\textcolor{tenue}{\$}}}1
           {á}{{\'a}}1 {é}{{\'e}}1 {í}{{\'i}}1 {ó}{{\'o}}1 {ú}{{\'u}}1 {ñ}{{\~n}}1,
}
\lstnewenvironment{terminal}{\lstset{style=terminal}}{}

%%%%%%%%%%%%%%%%%%%%%%%%%%%% Bibliografía %%%%%%%%%%%%%%%%%%%%%%%%%%%%%%%%%%%%%%%%

% Estilo IEEE (numérico, con las citas seguidas agrupadas: [1]-[4]) y los textos
% de la bibliografía en español
\usepackage[style=ieee,citestyle=ieee-comp,backend=biber]{biblatex}
\addbibresource{bibliografia.bib}
\renewcommand*{\bibfont}{\small}
\stdpunctuation % la coma tras las comillas, como en español («…», y no «…,»)
\DefineBibliographyStrings{spanish}{
  artno   = {art\adddot\addnbspace n\sptext{o}},
  jourser = {col\adddot},
  url     = {[En línea]\adddot\addspace Disponible en},
  urlseen = {Consultado el},
}

%%%%%%%%%%%%%%%%%%%%%%%%%%%% Siglas y acrónimos %%%%%%%%%%%%%%%%%%%%%%%%%%%%%%%%%%

% La primera vez que se usa una sigla se escribe completa; después, solo la sigla
\usepackage{acro}
\acsetup{
  make-links = true,
  list/name = {Siglas y acrónimos},
  list/heading = capitulo-en-indice,
}
\NewAcroTemplate[heading]{capitulo-en-indice}{%
  \chapter*{\acrolistname}%
  \addcontentsline{toc}{chapter}{\acrolistname}%
  \markboth{\MakeUppercase{\acrolistname}}{\MakeUppercase{\acrolistname}}}
% Siglas y acrónimos de la memoria (paquete acro). En el texto se usa \ac{clave}.
% El desarrollo corresponde letra a letra con la sigla: si es inglesa, en inglés
% y en cursiva; si es española, en español (README.md, «Cómo se escribe»).
% Formato de cada entrada:
%   \DeclareAcronym{crnn}{
%     short = CRNN,
%     long  = \textit{convolutional recurrent neural network},
%   }

\DeclareAcronym{bst}{
  short = BST,
  long  = \textit{Badminton Stroke-type Transformer},
}
\DeclareAcronym{crnn}{
  short = CRNN,
  long  = \textit{convolutional recurrent neural network},
}
\DeclareAcronym{fft}{
  short = FFT,
  long  = \textit{fast Fourier transform},
}
\DeclareAcronym{gru}{
  short = GRU,
  long  = \textit{gated recurrent unit},
}


%%%%%%%%%%%%%%%%%%%%%%%%%%%% Cabeceras y pies de página %%%%%%%%%%%%%%%%%%%%%%%%%%

\usepackage{fancyhdr}
\renewcommand{\chaptermark}[1]{\markboth{\chaptername\ \thechapter.\ #1}{}}
\renewcommand{\sectionmark}[1]{\markright{\thesection\ #1}}
\fancypagestyle{fancy}{%
  \fancyhf{}%
  \fancyhead[L]{\small\itshape\leftmark}%
  \fancyhead[R]{\small\rightmark}%
  \fancyfoot[C]{\small\thepage}%
  \renewcommand{\headrulewidth}{0.4pt}%
  \renewcommand{\footrulewidth}{0pt}}
\fancypagestyle{plain}{%
  \fancyhf{}%
  \fancyfoot[C]{\small\thepage}%
  \renewcommand{\headrulewidth}{0pt}%
  \renewcommand{\footrulewidth}{0pt}}

%%%%%%%%%%%%%%%%%%%%%%%%%%%% Datos fijos %%%%%%%%%%%%%%%%%%%%%%%%%%%%%%%%%%%%%%%%%

\newcommand{\escuela}{Escuela Técnica Superior\\de Ingeniería Informática}
\newcommand{\escuelalargo}{Escuela Técnica Superior de Ingeniería Informática}
\newcommand{\universidad}{Universidad Rey Juan Carlos}
\newcommand{\logoUniversidad}{config/logos/logoURJC.pdf}

%%%%%%%%%%%%%%%%%%%%%%%%%%%% Comentarios de revisión %%%%%%%%%%%%%%%%%%%%%%%%%%%%%

% \tutor{}, \cotutor{} y \alumno{} solo se ven con \drafttrue
\newif\ifdraft
\drafttrue
% Comentarios de revisión: se ven en el borrador (\drafttrue) y desaparecen en
% la versión final (\draftfalse)
\ifdraft
  \usepackage{todonotes}
\else
  \usepackage[disable]{todonotes}
\fi

\newcommand{\tutor}[1]{\todo[color=red!40, inline]{\footnotesize \textbf{Tutor:} #1}}
\newcommand{\cotutor}[1]{\todo[color=green!40, inline]{\footnotesize \textbf{Cotutor:} #1}}
\newcommand{\alumno}[1]{\todo[color=blue!40, inline]{\footnotesize \textbf{Alumno:} #1}}


%%%%%%%%%%%%%%%%%%%%%%%%%%%% Enlaces y referencias cruzadas %%%%%%%%%%%%%%%%%%%%%%%

% Enlaces sin color, para que la versión impresa salga en negro
\usepackage{hyperref}
\hypersetup{
  hidelinks,
  bookmarksnumbered,
  pdftitle  = {\titulotrabajo},
  pdfauthor = {\nombreautor},
  pdflang   = {es-ES},
}

% \cref{fig:x} escribe «figura 2.4» con el enlace; \Cref, a principio de frase
\usepackage[spanish]{cleveref}
\crefname{table}{tabla}{tablas}
\Crefname{table}{Tabla}{Tablas}
\crefname{appendix}{anexo}{anexos}
\Crefname{appendix}{Anexo}{Anexos}
\crefname{section}{sección}{secciones}
\Crefname{section}{Sección}{Secciones}
\crefname{subsection}{sección}{secciones}
\Crefname{subsection}{Sección}{Secciones}
